%Input preamble
%Style
\documentclass[12pt]{article}
\usepackage[top=1in, bottom=1in, left=1in, right=1in]{geometry}
\parindent 22pt
\usepackage{fancyhdr}

%Packages
\usepackage{adjustbox}
\usepackage{amsmath}
\usepackage{amsfonts}
\usepackage{amssymb}
\usepackage{bm}
\usepackage[table]{xcolor}
\usepackage{tabu}
\usepackage{color,soul}
\usepackage{makecell}
\usepackage{longtable}
\usepackage{multirow}
\usepackage[normalem]{ulem}
\usepackage{etoolbox}
\usepackage{graphicx}
\usepackage{tabularx}
\usepackage{ragged2e}
\usepackage{booktabs}
\usepackage{caption}
\usepackage{fixltx2e}
\usepackage[para, flushleft]{threeparttablex}
\usepackage[capposition=top,objectset=centering]{floatrow}
\usepackage{subcaption}
\usepackage{pdfpages}
\usepackage{pdflscape}
\usepackage{natbib}
\usepackage{bibunits}
\definecolor{maroon}{HTML}{990012}
\usepackage[colorlinks=true,linkcolor=maroon,citecolor=maroon,urlcolor=maroon,anchorcolor=maroon]{hyperref}
\usepackage{marvosym}
\usepackage{makeidx}
\usepackage{tikz}
\usetikzlibrary{shapes}
\usepackage{setspace}
\usepackage{enumerate}
\usepackage{rotating}
\usepackage{tocloft}
\usepackage{epstopdf}
\usepackage[titletoc]{appendix}
\usepackage{framed}
\usepackage{comment}
\usepackage{xr}
\usepackage{titlesec}
\usepackage{footnote}
\usepackage{longtable}
\newlength{\tablewidth}
\setlength{\tablewidth}{9.3in}
\setcounter{secnumdepth}{4}

\titleformat{\paragraph}
{\normalfont\normalsize\bfseries}{\theparagraph}{1em}{}
\titlespacing*{\paragraph}
{0pt}{3.25ex plus 1ex minus .2ex}{1.5ex plus .2ex}
\makeatletter
\pretocmd\start@align
{%
  \let\everycr\CT@everycr
  \CT@start
}{}{}
\apptocmd{\endalign}{\CT@end}{}{}
\makeatother
%Watermark
\usepackage[printwatermark]{xwatermark}
\usepackage{lipsum}
\definecolor{lightgray}{RGB}{220,220,220}
%\newwatermark[allpages,color=lightgray,angle=45,scale=3,xpos=0,ypos=0]{Preliminary Draft}

%Further subsection level
\usepackage{titlesec}
\setcounter{secnumdepth}{4}
\titleformat{\paragraph}
{\normalfont\normalsize\bfseries}{\theparagraph}{1em}{}
\titlespacing*{\paragraph}
{0pt}{3.25ex plus 1ex minus .2ex}{1.5ex plus .2ex}

\setcounter{secnumdepth}{5}
\titleformat{\subparagraph}
{\normalfont\normalsize\bfseries}{\thesubparagraph}{1em}{}
\titlespacing*{\subparagraph}
{0pt}{3.25ex plus 1ex minus .2ex}{1.5ex plus .2ex}

%Functions
\DeclareMathOperator{\cov}{Cov}
\DeclareMathOperator{\corr}{Corr}
\DeclareMathOperator{\var}{Var}
\DeclareMathOperator{\plim}{plim}
\DeclareMathOperator*{\argmin}{arg\,min}
\DeclareMathOperator*{\argmax}{arg\,max}

%Math Environments
\newtheorem{theorem}{Theorem}
\newtheorem{claim}{Claim}
\newtheorem{condition}{Condition}
\renewcommand\thecondition{C--\arabic{condition}}
\newtheorem{algorithm}{Algorithm}
\newtheorem{assumption}{Assumption}
\renewcommand\theassumption{A--\arabic{assumption}}
\newtheorem{remark}{Remark}
\renewcommand\theremark{R--\arabic{remark}}
\newtheorem{definition}[theorem]{Definition}
\newtheorem{hypothesis}[theorem]{Hypothesis}
\newtheorem{property}[theorem]{Property}
\newtheorem{example}[theorem]{Example}
\newtheorem{result}[theorem]{Result}
\newenvironment{proof}{\textbf{Proof:}}{$\bullet$}

%Commands
\newcommand\independent{\protect\mathpalette{\protect\independenT}{\perp}}
\def\independenT#1#2{\mathrel{\rlap{$#1#2$}\mkern2mu{#1#2}}}
\newcommand{\overbar}[1]{\mkern 1.5mu\overline{\mkern-1.5mu#1\mkern-1.5mu}\mkern 1.5mu}
\newcommand{\equald}{\ensuremath{\overset{d}{=}}}
\captionsetup[table]{skip=10pt}
%\makeindex

\setlength\parindent{20pt}
\setlength{\parskip}{0pt}

\newcolumntype{L}[1]{>{\raggedright\let\newline\\\arraybackslash\hspace{0pt}}m{#1}}
\newcolumntype{C}[1]{>{\centering\let\newline\\\arraybackslash\hspace{0pt}}m{#1}}
\newcolumntype{R}[1]{>{\raggedleft\let\newline\\\arraybackslash\hspace{0pt}}m{#1}}



%Logo
%\AddToShipoutPictureBG{%
%  \AtPageUpperLeft{\raisebox{-\height}{\includegraphics[width=1.5cm]{uchicago.png}}}
%}

\newcolumntype{L}[1]{>{\raggedright\let\newline\\\arraybackslash\hspace{0pt}}m{#1}}
\newcolumntype{C}[1]{>{\centering\let\newline\\\arraybackslash\hspace{0pt}}m{#1}}
\newcolumntype{R}[1]{>{\raggedleft\let\newline\\\arraybackslash\hspace{0pt}}m{#1}}

\newcommand{\mr}{\multirow}
\newcommand{\mc}{\multicolumn}

%\newcommand{\comment}[1]{}

%Other parameters
\newcommand{\noutcomes}{95}
\newcommand{\noutcomesexpp}{357}
\newcommand{\noutcomesexpm}{343}
\newcommand{\noutcomesexpf}{355}
\newcommand{\treatsubsabc}{$75\%$}
\newcommand{\treatsubscarec}{$74\%$}
\newcommand{\treatsubscaref}{$63\%$}

%Counts
%Males
\newcommand{\positivem}{$78\%$}
\newcommand{\positivesm}{$29\%$}

%Females
\newcommand{\positivef}{$78\%$}
\newcommand{\positivesf}{$31\%$}

%Counts, control substitution
%Males
\newcommand{\positivecsnm}{$47\%$}
\newcommand{\positivescsnm}{$15\%$}

\newcommand{\positivecsam}{$79\%$}
\newcommand{\positivescsam}{$29\%$}

%Females
%% no alternative
\newcommand{\positivecsnf}{$84\%$}
\newcommand{\positivescsnf}{$55\%$}

%% alternative
\newcommand{\positivecsaf}{$79\%$}
\newcommand{\positivescsaf}{$33\%$}

%Pooled

%Effects
%Males

%Females
\newcommand{\empf}{$8$}
\newcommand{\yearsedf}{$1.7$}



%Pooled

%CBA
%IRR
%Males
\newcommand{\irrm}{$15\%$}
\newcommand{\irrsem}{$5\%$}

%Females
\newcommand{\irrf}{$9\%$}
\newcommand{\irrsef}{$7\%$}

%Pooled
\newcommand{\irrp}{$13\%$}
\newcommand{\irrsep}{$5\%$}

%BC
%Males
\newcommand{\bcm}{$11.24$}
\newcommand{\bcsem}{$4.60$}

%Females
\newcommand{\bcf}{$2.35$}
\newcommand{\bcsef}{$1.09$}

%Pooled
\newcommand{\bcp}{$5.63$}
\newcommand{\bcsep}{$2.15$}

%NPV streams
%Pooled
\newcommand{\parincomenpvp}{$\$119,346$}

\newcommand*\leftright[2]{%
  \leavevmode
  \rlap{#1}%
  \hspace{0.5\linewidth}%
  #2}

\newcommand{\orth}{\ensuremath{\perp\!\!\!\perp}}%
\newcommand{\indep}{\orth}%
\newcommand{\notorth}{\ensuremath{\perp\!\!\!\!\!\!\diagup\!\!\!\!\!\!\perp}}%
\newcommand{\notindep}{\notorth}

\externaldocument{abc_treatmenteffects_appendix}

\begin{document}

\begin{comment}
\noindent Jorge - story line unclear still! (We will soon be booted out) \textbf{[JLG: That would be sad, but we simply seem not to be understanding each other, or maybe I'm just incompetent. But I keep explaining the same thing and you keep asking the same questions. Maybe these samples are much more limited than expected, even when we have done so much to them!]}\\

\noindent I am still looking for a clean story. At baseline, who is more disadvantaged boys or girls? In terms of:\\

\noindent {\bf [JLG: These are the numbers comparing all boys to all girls]}

\begin{itemize}
\item Education of mother
	\begin{itemize}
		\item Boys: 10.44 (s.e. .19)
		\item Girls: 10.38 (s.e. .21)
		\item Difference: 0.05 (s.e. .20)
\noindent \textbf{[JJH: The same.]} 
	\end{itemize}
\item Education of father
	\begin{itemize}
		\item Boys: 10.94 (s.e. 0.24)
		\item Girls: 10.92 (s.e. 0.16)
		\item Difference: 0.5 (s.e. .29)
\noindent \textbf{[JJH: The same.]} 
	\end{itemize}
\item Intact families 
	\begin{itemize}
		\item Boys: .28 (s.e. 0.05)
		\item Girls: .24 (s.e. 0.04)
		\item Difference: 0.04 (s.e. 0.07)
\noindent \textbf{[JJH: Slightly better.]} 
	\end{itemize}
\item Income + earnings of mother
	\begin{itemize}
		\item  Boys: 7,544 (s.e. 2,312) 
		\item  Girls: 5,158 (s.e. 1,531)
		\item Difference: 2,385 (s.e. 3,133) 
	\end{itemize}
\item Income + earnings of father
\noindent {\bf [JLG: not observed at baseline, observed later at some ages.] [JJH: Ok - so it's not baseline.] [JLG: Correct.]}
\end{itemize}

\noindent  {\bf [JLG: So, as we have discussed multiple times, boys seem to have advantage relative to girls if we take the differences in the levels of these variables seriously, but the differences are not statistically significant. That is why we use the index in Table 5. You keep saying in the text that the index in Table 5 is not at baseline, but as the note of the table clearly states, all of the measures are at baseline. You also say that the measure is arbitrary. That measure is a factor estimated from the following variables: Mother's age, mother's education, mother's IQ, mother's marital status, mother's employment, and number of siblings.] [JJH: But conceptually we want earnings at home including father at baseline.] [JLG: When we use the parental income aggregate in the paper (e.g., Tables 1 and 4) it is composed of what we have at that age. It includes mother, father, or parental income depending on the child's family situation. But recall that father income by itself is only observed at ages 12 and 21 and has very few observations.]}\\

\noindent {\bf [JLG: We estimate a factor using these variables to form the index, so the weight given to each of them is not arbitrary. Actually, the weights given in the HRI index are arbitrary given that it was set by the researchers at UNC.] [JJH: What do you mean by arbitrary?] [JLG: They constructed the index weighting the different components linearly according to some weights they literally made up.]}\\

\noindent {\bf [JLG: If you look at row [1], you will see that we reject the null of this index being the same across genders using the Rosenbaum test. The average level of this index is greater for boys than it is for girls, so we conclude that boys are more advantaged than girls.] [JJH: Do you form the factor using the same weights for boys and girls] [JLG: Yes, but we decided to do so because we failed to reject the null of equality.]}\\

\noindent {\bf [JLG: You also comment that ``everyone is disadvantaged in this sample!'' Of course, but we are obviously using the term disadvantaged as relatively disadvantaged within the experiment, not relative to the population of the US.]\\}

\noindent In the experiment, same question. What is the experimentally induced impact on family environment?

\begin{itemize} 
\item More fathers working? 
\item More mothers working? 
\item Family income? 
\end{itemize}

\noindent {\bf [JLG: An item-by-item breakout is going to leave to similar results as above and it could be arbitrary. But take a look at Table 4 and there are two rows relevant to your question.] [JJH: You misuse the word ``arbitrary.'' It's a legitimate question - more mothers work at baseline? Or after treatment. I thought the latter.] [JLG: Yes. It is the latter.] [JJH: Does it differ for girls.] [JLG: No. But we iterated this by email, and I think that we agreed on this. Also, in Section 5, I am answering a few questions related to this. Please also see that.] 

\begin{itemize} 
\item Parental Income: this is the set of variables containing mother's income, father's income (when we observe it), mother's and father's labor supply. As you see, the impact is REALLY large in terms of combining functions and the Rosenbaum tests indicate rejection. 
\item Parenting: this is basically the time series of the HOME scores from age 1 to age 8. Similarly, the effect of the experiment is large.
\end{itemize}

\noindent The sharp gender differences are there in the sense that in each of the outcome categories, girls benefit more from treatment than boys. Now, these are of course the treatment effects. But revisit the female-male gaps, pre- and post-treatment in Table 1 for each of the categories in Table 1. The differences are sharp. Most are significant!]}\\

\noindent I really do not see sharp differences in after or before effects. Before we were emphasizing greater vulnerability of boys:\\

\begin{itemize}
\item How do we know that this is not the story?
\item How to explain to referees?
\item Why now both? 
\item What is the strength of the evidence on either? 
\end{itemize}

\noindent {\bf [JLG: Yes, we were previously emphasizing greater vulnerability of boys. But we did not distill differences in baseline characteristics by gender to explain the differences in the treatment effects by gender to explain why the treatment effects differed depending on the counterfactual. Referee 3, the referee whose comments you praised, asked us to clarify that. That's why we explored the baseline differences in the control-group children who stayed at home and the control-group children who attended alternatives.\\ 

\noindent In Table 5, and please note that Table 5 is at baseline, we see sharp differences by gender by take-up of alternatives. \\

\noindent Within the control-group boys, the relatively more disadvantaged are sent more to alternatives. The relatively less disadvantaged are at home. Therefore, the treatment effect is greater for boys who went to alternatives - their counterparts who stayed at home were in better environments. Instead, these children were going to preschools of lower quality and were beginning the preschools at age 3 instead of at 8 weeks like ABC/CARE). \\

\noindent Now, for girls the opposite pattern appears, but not as strong. As a consequence, girls benefit more from treatment if compared to staying at home but the difference is not the same as in the case for boys.\\ 

\noindent Before, we explained this using quality. A careful inspection of the data implies that baseline characteristics could be an explanation as well. \\ 

\noindent Now, we discussed that it could very well be that both stories are there. But there's one big difference. Our information on quality is not good, while we have measures of disadvantage at baseline with which we can perform tests! So they both could be playing a role. As we discussed in person, it's going to be impossible to separate the stories in our sample or tell what's the strength for either.]}
\end{comment}

\pagebreak
\doublespacing

\pagebreak

\begin{comment}
\begin{center}
\noindent {\centering \textbf{Memo: June 13, 2018}}
\end{center}


\noindent \textbf{[JJH: JORGE and Anna: I am becoming deeply frustrated by this process as I am sure both of you are I cannot get a coherent story.]} \textbf{[JLG: We are providing answers to your questions day and night based on the data. We are not frustrated. We are putting hours into clarifying this. You seem insulted from the work that we do. You said two nights ago: ``SO WHY is an expanding and clarifying story an endless loop.'' Please give us some credit. We are grateful for your comments, but let's try to bounce ideas. Didn't you teach me that careful work takes time and pain? Our scientific responsibility is to tell you what the data looks like, not to make things easy for us and make up a story that you would like. Some would do so and we never will.]}\\

\noindent \textbf{[JJH: You say the treatment effect for family income is larger for BOYS.] [JLG: No. Baseline resources are greater but not the treatment effects.]}\\ 

\noindent \textbf{[JJH: Does this ``treatment effect'' include the higher baseline family income for boys?] [JLG: No. That would convolute pre- and post-treatment and none of the results in Table 1 or 4 do so. Why would they?]}\\

\noindent \textbf{[JJH: That's not a treatment effect.] [JLG: Thanks for the clarification, but it \textit{does not} include baseline income.} \\

\noindent \textbf{[JJH: As far as I can tell Boys start with higher income than girls (baseline).] [JLG: Correct.]}\\

\noindent \textbf{[JJH: Then the  mother is induced to work by program The INCREASE IN INCOME is greater for Girls than Boys.] [JLG: Yes.]}\\ 

\noindent \textbf{[JJH: That's the only logical story Boys start better Girls gain ground.] [JLG: Yes. That's why the gaps decrease.]}\\

\noindent \textbf{[JJH: The GROWTH in income is for girls.] [JLG: Yes.]}\\

\noindent \textbf{[JJH: The level of income may still be lower for treatment girls but they catch up.] [JLG: Yes.]}\\ 

\noindent \textbf{[JJH: The home scores go up for girls because of greater resources.] [JLG: Yes.]}\\

\noindent \textbf{[JJH: Fathers  are more likely to be present  for boys but the program does not induce in more fathers with boys or girls.] [JLG: Yes.]}\\

\noindent \textbf{[JJH: Fathers add income but fathers income does not increase.] [JLG: Yes.]}\\

\noindent \textbf{[JJH: I just keep getting part of the story from you.] [JLG: Is this so true? Or we are trying to understand the data and this is a difficult problem? Also, given the questions that you are bringing up, you seem to know the story.]}\\

\noindent \textbf{[JJH: I wont sign off unless we have a coherent story.] [JLG: We haven't asked you to sign off. We want to have a good study not a rotten publication.]}\\ 

\noindent \textbf{[JJH: As I said earlier the old draft from last month section 5 was frighteningly bad and has scared the hell out of me in regards to this paper.]} \textbf{[JLG: Is it constructive to bring this up now? As I said before, I am not sure why our work becomes so insulting to you. We worked and gave you a draft. You didn't like Section 5 and, as far as I am concerned, there's nothing bad about that. But we are coauthors and we can build a better Section 5 together. If the point is that we should take responsibility for writing a bad Section 5, then I am very happy to take it and move on. I thought writing a better Section 5 was part of the ``expanding and clarifying story,'' not that a section written months ago was going to disable us from moving on with the paper.]}\\

\noindent \textbf{[JJH: I feel I am fighting brick by brick house by house I still do not know if th treatment girlsget more incomefrm motherworking orfatherarriving do you know.]} \textbf{[JLG/AZ: ABC/CARE does not affect fathers presence. It does induce mothers to work. If the mothers are working and the father is present, then we do not know if the effect of that interaction given that we cannot distinguish father's income from mother's income. We do know that father's presence and family income are positively correlated. Please note that the treatment effects on parental income are actually larger for boys than for girls.]}\\

\noindent \textbf{[JJH: Table 4 is very inconclusive The description does not match the reality of the table  There is no Dominance of girls.]} \textbf{[JLG/AZ:  In Table 4, 7 out of 10 of the average effect sizes are larger for girls than for boys. In 9 out of 10 of the proportion of positive treatment effects, girls at least weakly dominate boys. In 8 out of 10 of the proportion of significant, positive treatment effects, girls at least weakly dominate boys. In addition, Table 1 shows that ABC/CARE reverses male-female gaps indicating that the treatment effects for girls dominate those for boys. This seems to indicate dominance of girls.]}\\

\noindent \textbf{[JJH: Figure 3 IS NOT at all clear and few differnces precisely determined I dropped 3.] [JLG: You requested that we brought that figure from the old draft the weekend before last because it ``told a story.'' I agree that this figure is not clear.]} \\

\noindent \textbf{[JJH: Its really unfortunate that we do not have POST TRatment Family enviornments in terms of income.]} \\\textbf{[JLG: Yes it is.]}\\

\noindent \textbf{[JJH: More mothers work  For gorls more motherswor but home score goes down  on net worse off.]}  \textbf{[JLG/AZ: Yes, treatment induces mothers to work. For girls, the home scores actually go up. This is seen in Table 4 under "Parenting," in which girls dominate boys. For example, the average effect size is 0.274 (and statistically significant) for girls, but 0.060 for boys. Recall that HOME scores include a subscale that negatively rates punishment, and the boys preform poorly in this specific subscale.]}\\

\noindent \textbf{[JJH: So its impossible to see if more fathers came home once child I nteatemt?]} \textbf{[JLG/AZ:It is possible to see if fathers are present. It is impossible to see their individual income. As we show in Tables D.57 and D.41 in the Appendix, treatment does not induce fathers to come back. For both boys and girls, the effect on father's presence after baseline is negative comparing to the control-group boys who stay at home. This is not consistent across ages though, indicating that fathers come and go in all subgroups. We verified in the data that indeed this is what is happening. Although we cannot formally disentangle the mediators of treatment, it seems that father's presence is not one of them.]}\\

\noindent \textbf{[JJH: Jorge, at the risk of facing another explosion, let me repeat my question which I have asked repeatedly and have not had answers.] [JLG: Please see the answers below.]}\\ 

\noindent \textbf{[JJH: Are family environments (other than HOME scores) changed by treatment?] [JLG: Yes. Apart from HOME scores, the other measure of family environment that we have is income. Both HOME scores and family income are positively affected by ABC/CARE.]} \\

\noindent \textbf{[JJH: Is part of the better female treatment effects due to better environments at home for girls?] [JLG: Yes. Improved environments. Part of the better female treatment effects is due to the worse environments at home for girls than boys. The control-group boys are more advantaged than the control-group girls. This is consistent with there being more positive treatment effects for females.] [JJH: I keep asking the same question - Better environments for treatment-group girls?] [JLG: Yes. Read the answers. It says ``Yes. Improved environments.'']}\\

\noindent \textbf{[JJH: More mothers induced to work for girls?] [JLG: Yes. A larger percentage of control boys have working mothers at baseline, so the treatment effects on maternal employment are generally larger for girls than for boys.] [JJH: Ah - an answer.]} \\

\noindent \textbf{[JJH: more family income?] [JLG: Yes. The treatment results in more family income. This is shown in Tables 1 and 4. The effect on family income is larger for boys than for girls. Remember that more boys have fathers at home than girls. The fact that the treatment effect on maternal employment is larger for girls is not inconsistent with the treatment effect on family income being larger for boys.] [JJH: Male income income up and female employment up?] [JLG: No. This is based on mother. Father's data on income is garbage. It was less than 20 observations.]} \\

\noindent \textbf{[JJH: More fathers coming back to an improved environment?] [JLG: No. Treatment does not induce fathers to ``come back,'' although families with fathers present are more advantaged along other measures, e.g., family income.] [JJH: OK.]} \\

\noindent \textbf{[JJH: You said we did not know father's income at baseline When is father'��s income measured?] [JLG: It is measured at ages 12 and 21, although there are very few observations (28 and 4, respectively).] [JJH: Use this! Are fathers working more.] [JLG: We do now know if they work more. And these are 28 and 4 observations. How can we use 4 observations? Even 28 is very very bad. We have 8 controls. This would not be serious and we cannot detect anything when we do it anyway.]} \\

\noindent \textbf{[JJH: Do we know if father present at baseline? Yes or so it seems] [JLG: Yes. Whenever the paper refers to father's presence, it is at baseline.]}\\ 

\noindent \textbf{[JJH: Do we know if fathers increase their presence at home if the son is in treatment?] [JLG: Fathers do not increase their presence at home if the son is in treatment. This is evident in Figure 1 of the old, submitted draft. This is attached.]}\\

\noindent \textbf{[JJH: Do we know this question for daughters?] [JLG: Fathers do not increase their presence at home if the daughter is in treatment. It is actually the case that more fathers are home for control girls than treatment girls. This difference is only statistically significant at age 4.5 when 17\% of the treatment girls have fathers present (observed for 30 subjects) and 31\% of the control girls have fathers present (observed for 36 subjects). Please note that at this age, the number of observations of fathers present is very small.]}\\

\noindent \textbf{[JJH:  Do we know differential treatment effects?] [JLG: As you can see in the previous answer, we don't have great measures of disadvantage after baseline. However, as Tables 1 and 4 show, ABC/CARE differentially affects parental income and parenting by gender. The effects on parental income are larger for boys than for girls. The effects on parenting are larger for girls than for boys.]}\\ 

\noindent \textbf{[JJH: To what extent does treatment improve family incomes of treatment relative to control?] [JLG: This is displayed in Table 4 and the differential effect in Table 1.]}\\ 

\noindent \textbf{[JJH: More mothers are INDUCED to work above baseline correct?] [JLG: Yes. Treatment increases maternal employment by about 10\%. A larger percent of control boys have working mothers at baseline, so the treatment effects on maternal employment are generally larger for girls than for boys.]}\\

\noindent \textbf{[JJH: Is the greater family income of treatments from the mothers or the INCREASE in the presence of the father] [JLG:  When fathers stay at home, the family income increases. Unfortunately, we do not have good documentation of who this family income is coming from, so it is hard to determine the exact effect of the father. It is possible that the father provides additional income. It is also possible the father provides childcare (or additional relatives who can provide childcare) allowing the mothers to work more or in better jobs.]}\\

\noindent \textbf{[JLG: We do not have gender of other siblings. Ideally, we would have a time series of the measure of disadvantage for each child. I think we both agree on that. But we don't have that. We do have: A time series of the home score, for which we do have results in the paper both in Table 1 and Table 4, which are by gender and comparing boys and girls. We call this "parenting" in the paper. Parental income:  which is an aggregate of mother's, father's and both mother and father income. The reason for using the different variables is that at different ages we observe different variables. For this we also present results by gender and comparing genders in Tables 1 and 4.]}

\end{comment}

\begin{center}
\noindent {\centering \textbf{Memo: June 14, 2018}}
\end{center}

\noindent \textbf{[JJH: We want to know Why the Treatment effects are bigger for girls.] [JLG: Yes. I agree and the evidence indicates, as you pointed out yesterday, that the environments are improving more for girls than for boys. This is in terms of growth.]}\\ 

\noindent \textbf{[JJH: We know at baseline Girls are in worse circumstances.] [JLG: Yes.]}\\

\noindent \textbf{[JJH: What we want to know is why they IMPROVE That involves analyses BEYOND baseline.] [JLG: Yes. Please see Table~\ref{table:massiveall} in this document. In terms of the environment of girls and boys we have the income and labor supply of their mother and the HOME scores (which include interaction with both the mother and the father, when the father is around). This is extracted from the paper, but may be clearer when considered separately. These two aggregates post-baseline, as described in the table. The reason why we aggregate this is because otherwise the estimates get very noisy.]}\\

\begin{sidewaystable}[H]
\begin{threeparttable}
\caption{Excerpt from Table 4: Combining Functions and Exact Non-Parametric Tests} \label{table:massiveall}
\centering
\begin{tabular}{l c c c c}
\toprule
 & Average & \% $ >0 $ & \% $ >0 $ , Significant & Rosenbaum \\
 & Effect Size & Treatment Effect & Treatment Effect & $ p $ -value \\
\midrule
\textbf{Parental Income} & \mc{4}{c}{[Parental income (1.5, 2.5, 3.5, 4.5) and mother works (1.5, 2.5, 3.5, 4.5)]} \\
 &      & &    & \\
\quad Females & \textbf{    0.283} & \textbf{  100.000} & \textbf{   75.000}  & \textbf{ 0.086} \\
\quad Males &  \textbf{    0.157} & \textbf{  100.000} & \textbf{   37.500}  & 0.147 \\
\midrule
\textbf{Parenting} & \mc{2}{c}{[HOME scores (0.5, 1.5, 2.5, 3.5, 4.5)]} & & \\
 &      & &    & \\
\quad Females &  \textbf{    0.274} & \textbf{  100.000} & \textbf{   80.000} & 0.602 \\
\quad Males &      0.060 & \textbf{   80.000} &     0.000 & 0.147 \\
\bottomrule
\end{tabular}
\begin{tablenotes}
\footnotesize
\item \textbf{Note:} This table displays summaries of treatment effects by outcome category and gender. Each of the panels contains statistics calculated using outcomes grouped by category. The average effect size is calculated by averaging over the effect size of the outcomes in the outcome category. The effect sizes of the individual outcomes are calculated by dividing the treatment-control mean difference by the standard deviation of the control group. We present bootstrapped $p$-values. For the proportion of outcomes that are positive and significant, we do a ``double bootstrap'' procedure. The null hypothesis for the average effect sizes is that they are 0. The null hypothesis for the proportion of outcomes that are (significantly) positive is that they are (10\%) 50\%. Bolded statistics are significant at the 10\% level. The Rosenbaum $p$-value originates from a test where the null is a common joint distribution of the variables in each category. For computational simplicity, we approximate the exact $p$-values using asymptotic $p$-values. Rosenbaum (2005) presents several simulation exercises showing that the validity of this approximation. Statistics significant at the $0.10$ level are bolded.
\end{tablenotes}
\end{threeparttable}
\end{sidewaystable}

\noindent \textbf{[JJH: You seem bound and determined to stay at baseline and that's non-responsive to the questions of the referees and to fundamental curiosity and scientific discourse.] [JLG: No. I want to show post-baseline, as indicated in the previous Table~\ref{table:massiveall}. Now, father's income comes next.]}\\

\noindent \textbf{[JJH: You are some saying 28 observations is too small It may be But please I understand samples may be small I understand you would like more So would I But please understand that a baseline only analysis is essentially unsatisfactory It does not  explain CHANGE except possibly through an interaction.] [JLG: Yes. But I think we should construct that story based on Table~\ref{table:massiveall} in this document and not on father's income, for which you can see the results next in Table~\ref{table:father} in this document. Now, one could ask why not add father income into the aggregates in Table~\ref{table:massiveall} in this document. The reason is that if we do so the number of observations drops so much and we are back in trouble.]}\\ 

\begin{table}[H]
\begin{threeparttable}
\caption{Description of Father's Income} \label{table:father}
\centering
\begin{tabular}{l c c c c c c}
\toprule
& Age & \mc{2}{c}{$N$} & \mc{2}{c}{Mean} & $\Delta$ \\
\cmidrule(lr){3-4} \cmidrule(lr){5-6}
& 	& Control & Treatment & Control & Treatment & \\
\midrule
Males & 12 & 8 & 8 & $25,466.05$ & $25,147.62$ & -318.43 \\
 & 21 & 0 & 4 & --- & $11,507.50$ & --- \\
 \midrule
Females & 12 & 7 & 5 & $26,466.40$ & $52,612.00$ & $26,145.60$ \\
 & 21 & 0 & 0 & --- & --- & --- \\
\bottomrule
\end{tabular}
\begin{tablenotes}
\footnotesize
\item \textbf{Note:} This table describes father's income by treatment status, displaying age of measurement, number of observations in each treatment regime, control and treatment means, and average treatment effect. Both variables are only available for subjects in ABC (not CARE). All values are in 2014 USD.
\end{tablenotes}
\end{threeparttable}
\end{table}

\noindent \textbf{[JJH: You seem to think I do not understand that you have analyzed baseline I do understand very well now although the writing in previous drafts was exceedingly unclear.] [JLG: No. I am just trying to clarify things in response to the questions, but I mean no disrespect.]}\\ 

\noindent \textbf{[JJH: But now we need more however imperfect.] [JLG: I think that the parental income and HOME scores tell the story.]}\\ 

\noindent \textbf{[JJH: The story in 5 is weak and any fragments will help Read what I wrote carefully.]}












 






\end{document}